\section{Calibration Knobs for vLLM or SGLang}

This section summarizes which simulator parameters should be adjusted when the
goal is to match measured request-level behavior from a real serving system such
as vLLM or SGLang. The code paths that consume these knobs are primarily
\texttt{ServingScaleModel::run()} in
\texttt{astra-sim/analytical/ServingScaleModel.cc},
\texttt{ServingCostModel::estimate\_stage\_breakdown(...)} in
\texttt{astra-sim/workload/ServingCostModel.cc}, and the runtime schedulers in
\texttt{ColocatedServingRuntime.cc} and \texttt{PdServingRuntime.cc}.

\subsection{Calibration Objective}

Let \(\theta\) denote the set of simulator knobs selected for fitting. A useful
request-level calibration objective is:

\begin{equation*}
\begin{aligned}
\theta^\star
= \arg\min_{\theta} \\
& \sum_{r \in \mathcal{R}}
\Bigl[
w_{\mathrm{ttft}}
\bigl(
\mathrm{TTFT}_{\mathrm{sim}}(r; \theta) -
\mathrm{TTFT}_{\mathrm{meas}}(r)
\bigr)^2 \\
& {}+ w_{\mathrm{tpot}}
\bigl(
\mathrm{TPOT}_{\mathrm{sim}}(r; \theta) -
\mathrm{TPOT}_{\mathrm{meas}}(r)
\bigr)^2 \\
& {}+ w_{\mathrm{e2e}}
\bigl(
\mathrm{E2E}_{\mathrm{sim}}(r; \theta) -
\mathrm{E2E}_{\mathrm{meas}}(r)
\bigr)^2
\Bigr]
\end{aligned}
\end{equation*}

If the goal is SLO goodput rather than raw latency alignment, the objective can
be extended with a goodput term:

\begin{equation*}
\begin{aligned}
\theta^\star
= \arg\min_{\theta} \\
& \Bigl(
\mathcal{L}_{\mathrm{latency}}(\theta) \\
& {}+ w_{\mathrm{goodput}}
\bigl(
\mathrm{Goodput}_{\mathrm{sim}}(\theta) -
\mathrm{Goodput}_{\mathrm{meas}}
\bigr)^2
\Bigr)
\end{aligned}
\end{equation*}

\subsection{Parameters That Should Usually Be Fixed, Not Fit}

These fields should normally be copied from the real deployment rather than
treated as free parameters:

\begin{longtable}{p{0.28\textwidth} p{0.62\textwidth}}
\toprule
Config field family & Why it should usually be fixed from the real system \\
\midrule
\endhead
\texttt{topology.*\_layout.tp\_degree} &
Tensor-parallel degree is an explicit engine launch choice. \\
\texttt{topology.*\_layout.pp\_degree} &
Pipeline-parallel degree is a structural deployment choice. \\
\texttt{topology.*\_layout.ep\_degree} &
Expert-parallel degree is fixed by the MoE serving configuration. \\
\texttt{topology.*\_layout.dp\_attention\_degree} &
KV sharding ownership should follow the real decode-side layout. \\
\texttt{topology.*\_layout.dp\_replica\_count} &
Replica count should match the number of independent serving domains. \\
\texttt{cluster.devices[*]} &
Peak FLOPs, memory size, and memory bandwidth should reflect the real GPU SKU. \\
\texttt{interconnects[*]} &
Link bandwidth, latency, and efficiency should reflect NVLink, PCIe, or network fabric. \\
\texttt{model.*} &
Layer count, hidden size, KV heads, expert count, and byte widths should match the deployed model. \\
\texttt{runtime.architecture} &
Choose colocated or \texttt{pd\_disaggregated} to match the real engine architecture. \\
\bottomrule
\end{longtable}

\subsection{Primary Fit Knobs}

Once the structural deployment is fixed, the main calibration parameters are:

\begin{longtable}{p{0.30\textwidth} p{0.28\textwidth} p{0.32\textwidth}}
\toprule
Knob family & Config fields & What measured behavior it mainly controls \\
\midrule
\endhead
Global compute scaling &
\texttt{compute\_scale} &
Uniform up or down adjustment of all compute-side terms. \\
Global communication scaling &
\texttt{comm\_scale} &
Uniform up or down adjustment of TP, PP, EP, and DP-attention communication terms. \\
Prefill stage compute overrides &
\texttt{cost\_model.prefill.attention\_compute\_ns\_per\_token},
\texttt{cost\_model.prefill.ffn\_or\_expert\_compute\_ns\_per\_token} &
Single-request prefill TTFT and prompt-length scaling. \\
Decode stage compute overrides &
\texttt{cost\_model.decode.attention\_compute\_ns\_per\_token},
\texttt{cost\_model.decode.ffn\_or\_expert\_compute\_ns\_per\_token} &
Single-request TPOT and output-length scaling. \\
TP communication overrides &
\texttt{cost\_model.<stage>.tp\_collective\_ns\_per\_token} &
Sensitivity to tensor-parallel degree and collective cost. \\
PP communication overrides &
\texttt{cost\_model.<stage>.pp\_activation\_ns\_per\_token} &
Sensitivity to pipeline activation traffic and stage bubble cost. \\
EP communication overrides &
\texttt{cost\_model.<stage>.ep\_dispatch\_ns\_per\_token} &
MoE token dispatch and combine overhead. \\
DP-attention overrides &
\texttt{cost\_model.<stage>.dp\_attention\_sync\_ns\_per\_token} &
Decode-side KV synchronization overhead. \\
Batching efficiency &
\texttt{cost\_model.prefill.batch\_efficiency},
\texttt{cost\_model.decode.batch\_efficiency} &
How stage runtime scales with batch size rather than single-request cost. \\
Decode contention factor &
\texttt{cost\_model.decode.interference\_factor} &
How much decode slows under continuous batching or mixed load. \\
PD transfer model &
\texttt{pd.transfer.latency\_ns},
\texttt{pd.transfer.bandwidth\_bytes\_per\_s},
\texttt{pd.transfer.efficiency},
\texttt{pd.transfer.bytes\_per\_prompt\_token} &
Prefill-to-decode handoff overhead in disaggregated serving. \\
Scheduler concurrency caps &
\texttt{scheduler.max\_running\_requests},
\texttt{scheduler.max\_decode\_batch\_requests},
\texttt{scheduler.max\_prefill\_batch\_tokens},
\texttt{scheduler.prefill\_max\_requests} &
Queueing delay, batch shape, and admission pressure. \\
PD worker counts &
\texttt{pd.prefill\_workers},
\texttt{pd.decode\_workers} &
Stage-level concurrency and overlap in PD mode. \\
\bottomrule
\end{longtable}

\subsection{Direct Estimators for the Main Fit Knobs}

If stage measurements are available from profiling or tracing, the global scale
factors can be initialized from:

\begin{equation*}
s_c \approx
\frac{T_{\mathrm{compute,meas}}}
{\widehat{T}_{\mathrm{compute,model}}}
\end{equation*}

\begin{equation*}
s_m \approx
\frac{T_{\mathrm{comm,meas}}}
{\widehat{T}_{\mathrm{comm,model}}}
\end{equation*}

If the calibration is performed through direct per-token overrides, a stage
component coefficient can be initialized from:

\begin{equation*}
c_x^{\mathrm{stage}} \approx
\frac{T_{x,\mathrm{meas}}}
{m_{\mathrm{stage}} N_{\mathrm{tok}}}
\end{equation*}

For PD transfer, if \(Q_{\mathrm{KV}}\) is known and latency is measured
separately, a first estimate of effective bandwidth is:

\begin{equation*}
BW_{\mathrm{pd}} \epsilon_{\mathrm{pd}} \approx
\frac{Q_{\mathrm{KV}}}
{T_{\mathrm{transfer,meas}} - L_{\mathrm{pd}}}
\end{equation*}

If the transfer bytes are not directly observable, the override field can be
fit from:

\begin{equation*}
q_{\mathrm{kv,override}} \approx
\frac{Q_{\mathrm{KV,meas}}}{N_p}
\end{equation*}

\subsection{Recommended Calibration Order}

\begin{enumerate}
\item Fix \texttt{model}, \texttt{cluster}, \texttt{interconnects}, and
\texttt{topology} to match the real deployment. Do not use cost overrides to
hide a topology mismatch.
\item Fit \texttt{compute\_scale} and \texttt{comm\_scale} on clean,
single-request measurements with minimal queueing.
\item Fit prefill component overrides using prompt-only or TTFT-dominated runs:
\texttt{cost\_model.prefill.*}.
\item Fit decode component overrides using long-generation runs:
\texttt{cost\_model.decode.*}.
\item Fit \texttt{batch\_efficiency} and
\texttt{cost\_model.decode.interference\_factor} using concurrency sweeps and continuous
batching experiments.
\item If using PD, fit \texttt{pd.transfer.*} and then tune
\texttt{pd.prefill\_workers}, \texttt{pd.decode\_workers},
\texttt{topology.prefill\_layout}, and \texttt{topology.decode\_layout} to match overlap and
handoff behavior.
\item Finish with trace-level validation against TTFT, TPOT, E2E, and SLO
goodput on realistic request mixes.
\end{enumerate}

\subsection{vLLM-Specific Emphasis}

For colocated vLLM-style calibration, the highest-leverage knobs are usually:

\begin{itemize}
\item \texttt{compute\_scale}
\item \texttt{comm\_scale}
\item \texttt{cost\_model.decode.batch\_efficiency}
\item \texttt{cost\_model.decode.interference\_factor}
\item \texttt{scheduler.max\_decode\_batch\_requests}
\item \texttt{scheduler.max\_running\_requests}
\end{itemize}

The reason is that vLLM latency is often dominated by decode-side continuous
batching behavior, so the decode stage and admission caps usually matter more
than PD transfer parameters.

\subsection{SGLang-Specific Emphasis}

For SGLang-style PD or prefill or decode disaggregation, the highest-leverage
knobs are usually:

\begin{itemize}
\item \texttt{prefill\_layout}
\item \texttt{decode\_layout}
\item \texttt{pd.prefill\_workers}
\item \texttt{pd.decode\_workers}
\item \texttt{pd.transfer.latency\_ns}
\item \texttt{pd.transfer.bandwidth\_bytes\_per\_s}
\item \texttt{pd.transfer.efficiency}
\item \texttt{pd.transfer.bytes\_per\_prompt\_token}
\end{itemize}

The reason is that disaggregated serving is highly sensitive to handoff cost,
decode concurrency, and whether the prefill and decode tiers use the same or
different TP, PP, EP, and DP-attention layouts.

\subsection{What Should Be Changed First When the Fit Is Wrong}

If TTFT is wrong but TPOT is reasonable, change:

\begin{itemize}
\item \texttt{cost\_model.prefill.*}
\item \texttt{scheduler.max\_prefill\_batch\_tokens}
\item \texttt{scheduler.prefill\_max\_requests}
\item \texttt{pd.transfer.*}
\end{itemize}

If TPOT is wrong but TTFT is reasonable, change:

\begin{itemize}
\item \texttt{cost\_model.decode.*}
\item \texttt{cost\_model.decode.batch\_efficiency}
\item \texttt{cost\_model.decode.interference\_factor}
\item \texttt{scheduler.max\_decode\_batch\_requests}
\end{itemize}

If both TTFT and TPOT are wrong in the same direction, change:

\begin{itemize}
\item \texttt{compute\_scale}
\item \texttt{comm\_scale}
\item \texttt{cluster.devices[*]}
\item \texttt{interconnects[*]}
\end{itemize}

If latency percentiles are wrong mainly because queueing is wrong, change:

\begin{itemize}
\item \texttt{scheduler.max\_running\_requests}
\item \texttt{pd.prefill\_workers}
\item \texttt{pd.decode\_workers}
\item \texttt{topology.*\_layout.dp\_replica\_count}
\end{itemize}
